\section{Discussion}
\label{sec:discussion}

RPL-ClusterIDS shows that energy- and context-aware clustering can orchestrate hierarchical intrusion detection in RPL without modifying the routing protocol itself.
The design is most advantageous when attacks are spatially correlated, when application traffic is heterogeneous, and when the network cannot afford uniform high-intensity monitoring on every node.

\subsection{Why Clustering Matters}
The ablation design in Table~\ref{tab:ablation} isolates four architectural components: clustering, CH threshold verification, context adaptation, and energy-aware policy.
Clustering concentrates advanced analysis on a small set of capable heads rather than on every mote and bounds collaboration through LAS summarization and thresholded inter-head signaling.
Flat distributed IDS designs may approach similar detection rates on small networks, but their alert traffic typically grows much faster with node count; this hypothesis is testable via Fig.~\ref{fig:alert-overhead} once the scalability campaign completes.

\subsection{Positioning Relative to Baselines}
The centralized baseline (B1) benefits from global visibility at the border router but introduces a single point of failure and higher detection latency for spatially localized attacks.
RPL-ClusterIDS avoids this limitation through hierarchical distributed detection: members perform local screening while cluster heads aggregate evidence, with optional border-router corroboration for network-wide campaigns.
Additional baselines (flat distributed B2 and centralized-with-CH B3) are defined in the firmware and will be evaluated in the full campaign.

\subsection{Operating-Mode Trade-offs}
The three-mode policy (\textsc{Full}, \textsc{Balanced}, \textsc{Eco}) makes the security--lifetime compromise explicit and measurable.
Table~\ref{tab:modes} and Fig.~\ref{fig:mode-sensitivity} quantify this trade-off for the pilot campaign; full statistical validation is left for future work.
Practitioners can select a mode based on application criticality ($C0$--$C3$) and remaining network energy.

Limitations are discussed separately in Section~\ref{sec:limitations}.
